\section*{ID7218: Incorporating the Accumulated-Force Gravity Correction (2)}
\paragraph{Metadata.}
PiID: 3576164773794 | RTEID: RTE:P28029.0970:T2.2010 | UUID: 04fa0b2e-031f-5772-b158-215d836d00fb

\paragraph{Canonical Statement.}
\textbackslash{}[ \textbackslash{}delta T\_\{ij\} \textbackslash{}sim \textbackslash{}alpha \textbackslash{}frac\{G\textbackslash{},M\_\{\textbackslash{}rm acc\}\}\{R\textasciicircum{}3\} \textbackslash{}, \textbackslash{}hat n\_i \textbackslash{}hat n\_j, \textbackslash{}]

\paragraph{Canonical Equation.}
\[
\delta T_{ij} \sim \alpha \frac{G\,M_{\rm acc}}{R^3} \, \hat n_i \hat n_j,
\]

\paragraph{Dictionary.}
Incorporating the Accumulated-Force Gravity Correction (2) — δ T\_ij ∼ α fracG M\_rm accR\textasciicircum{}3 hat n\_i hat n\_j

\paragraph{Encyclopedia.}
Incorporating the Accumulated-Force Gravity Correction (2) is a scaling / order-of-magnitude relation, written δ T\_ij ∼ α fracG M\_rm accR\textasciicircum{}3 hat n\_i hat n\_j. It expresses a scaling or proportionality, capturing how the quantity grows with its parameters. It is built from δ, α, T\_ij, G, M\_ acc, R. Within the Trawin Zero Point registry it serves as one element of the framework's mathematical narrative.
